% Class-compatibility preamble retained for a reader excerpt of
%   Dave Witte Morris, "A Cohomological Proof that Real Representations of
%   Semisimple Lie Algebras Have Q-Forms", SIGMA 11 (2015) 034,
%   DOI 10.3842/SIGMA.2015.034, CC BY 4.0.
%
% The source pins the publisher class sigma, which is not part of the delivery
% baseline, so this file supplies the amsart baseline and then copies the author's
% own declaration block mechanically and unchanged:
%   - pinned-original-source.tex lines 11-31: the author's DeclareMathOperator
%     list, the equation numbering and the author's newtheorem declarations;
%   - pinned-original-source.tex lines 35-43, which the author writes immediately
%     after \begin{document} in the publisher layout: the author's own \Res,
%     \ResOne and \Lie macros and \allowdisplaybreaks.
% The publisher header commands of the sigma class (\FirstPageHeading,
% \ArticleName and the like) and the macro \PaperNumber they require are not
% reproduced, because this excerpt carries neither the journal front matter nor
% anything else that uses them.
% No mathematical macro, statement, symbol or assumption is introduced, changed
% or removed by this file.
\documentclass[10pt,a4paper]{amsart}
\usepackage{amssymb,amsmath,amsthm}
\usepackage{enumerate}
\usepackage{hyperref}

\DeclareMathOperator{\Cliff}{Clif\/f}
\DeclareMathOperator{\Gal}{Gal}
\DeclareMathOperator{\Inn}{Inn}
\DeclareMathOperator{\Mat}{Mat}
\DeclareMathOperator{\Out}{Out}
\numberwithin{equation}{section}

\newtheorem{prop}{Proposition}[section]
\newtheorem{cor}[prop]{Corollary}
\newtheorem{lem}[prop]{Lemma}
{\theoremstyle{definition}
\newtheorem{defn}[prop]{Definition}
\newtheorem{notation}[prop]{Notation}
\newtheorem{rem}[prop]{Remark}
\newtheorem{ques}[prop]{Question}
\newtheorem{eg}[prop]{Example}
\newtheorem{erratum}[prop]{Erratum}}

\newcommand{\Res}[1]{\mathop{\mathbf{R}_{#1}}}
\newcommand{\ResOne}[1]{\mathop{\mathbf{R}^{(1)}_{#1}}}
\newcommand{\Lie}[1]{\mathfrak{\lowercase{#1}}}
\allowdisplaybreaks
\usepackage{cleveref}
\crefname{originallemma}{original lemma}{original lemmas}
\Crefname{originallemma}{Original lemma}{Original lemmas}
\makeatletter
\newlabel{RunivIff}{{2.1}{source}{}{sourcepointer-1}{}}
\newlabel{RunivIff@cref}{{[originallemma][0][]2.1}{[1][0][]source}{}{}{}}
\makeatother
\crefname{originaldefinition}{original definition}{original definitions}
\Crefname{originaldefinition}{Original definition}{Original definitions}
\makeatletter
\newlabel{RunivIff-D}{{4}{source}{}{sourcepointer-2}{}}
\newlabel{RunivIff-D@cref}{{[originaldefinition][0][]4}{[1][0][]source}{}{}{}}
\makeatother
\crefname{originaldefinition}{original definition}{original definitions}
\Crefname{originaldefinition}{Original definition}{Original definitions}
\makeatletter
\newlabel{RunivIff-DivAlg}{{3}{source}{}{sourcepointer-3}{}}
\newlabel{RunivIff-DivAlg@cref}{{[originaldefinition][0][]3}{[1][0][]source}{}{}{}}
\makeatother
\crefname{originalremark}{original remark}{original remarks}
\Crefname{originalremark}{Original remark}{Original remarks}
\makeatletter
\newlabel{OnlyZStar}{{2.4}{source}{}{sourcepointer-4}{}}
\newlabel{OnlyZStar@cref}{{[originalremark][0][]2.4}{[1][0][]source}{}{}{}}
\makeatother
\title{An R-universal semisimple Lie algebra becomes inner over an imaginary quadratic extension with the prescribed star-action}
\author{Dave Witte Morris}
\date{doi:10.3842/SIGMA.2015.034; CC BY 4.0}
\begin{document}
\maketitle
\noindent\textit{Editorial note.} All delivered mathematical text is the author's own, copied from the pinned publisher source sources/algebra-fl02-a-cohomological-proof-that-real-re/pinned-original-source.tex; no statement or proof step was rewritten, reordered or removed.\par\medskip
\noindent\textit{Editorial note.} Preamble provenance: the source pins the publisher class sigma, which is not part of the delivery baseline, so this card uses sources/algebra-fl02-a-cohomological-proof-that-real-re/retained-preamble.tex, which supplies the amsart baseline, copies the author's DeclareMathOperator list, equation numbering and newtheorem declarations from pinned source lines 11-31, and copies the author's own \textbackslash{}Res, \textbackslash{}ResOne and \textbackslash{}Lie macros from pinned source lines 35-43, which the author writes immediately after \textbackslash{}begin\{document\}. No macro was invented and no mathematical symbol was changed.\par\medskip
\noindent\textit{Editorial note.} Bibliography provenance: the counted statement and its proof contain no citation, so this card carries no bibliography. The source has a thebibliography environment at pinned source lines 954-1016 whose three entries are used elsewhere in the paper; nothing from it is needed here and no entry was reconstructed.\par\medskip
\noindent\textit{Editorial note.} External references are navigation pointers only; every deliberately kept step of the author's argument stays inside the counted spans. Their printed numbers are the author's own numbering under the author's newtheorem declarations, computed mechanically by the lane's scanner.\par\medskip
\noindent\textit{Allowed standard prerequisites (editorial).}\begin{itemize}
\item The notion of a Lie algebra being inner over an extension L, and the star-action of Gal(L/Q) compared with the star-action of Gal(C/R). Set up in the author's definitions and notation before the counted statement (source lines 135-216). Fixes the two properties that the statement asserts about the extension L.
\item The Tits algebra D(Lie G, Q)(lambda) and the statement that its centre is the minimal extension over which the algebra becomes inner. The author's proposition labelled Beta=D at source lines 393-406, used inside the counted proof. Identifies the centre of the division algebra with the minimal splitting extension and lets the proof conclude that it is imaginary quadratic.
\end{itemize}\medskip
\ifdefined\hypertarget\hypertarget{sourcepointer-1}{}\fi
\noindent\textit{Source reference.} Original lemma 2.1: sources/algebra-fl02-a-cohomological-proof-that-real-re/pinned-original-source.tex, source lines 416--416. Navigation pointer only: the counted proof invokes the author's equivalence for R-universality in its last paragraph; the equivalence has its own card in this lane and is declared here as a named standard prerequisite.\par\medskip
\ifdefined\hypertarget\hypertarget{sourcepointer-2}{}\fi
\noindent\textit{Source reference.} Original definition 4: sources/algebra-fl02-a-cohomological-proof-that-real-re/pinned-original-source.tex, source lines 432--432. Navigation pointer only: it is the fourth labelled item of that equivalence, cited by the counted proof once.\par\medskip
\ifdefined\hypertarget\hypertarget{sourcepointer-3}{}\fi
\noindent\textit{Source reference.} Original definition 3: sources/algebra-fl02-a-cohomological-proof-that-real-re/pinned-original-source.tex, source lines 428--428. Navigation pointer only: it is the third labelled item of that equivalence, cited by the counted proof once.\par\medskip
\ifdefined\hypertarget\hypertarget{sourcepointer-4}{}\fi
\noindent\textit{Source reference.} Original remark 2.4: sources/algebra-fl02-a-cohomological-proof-that-real-re/pinned-original-source.tex, source lines 462--462. Navigation pointer only: the counted proof invokes this remark on the root-lattice case in its closing argument; its content is declared as a named standard prerequisite.\par\medskip
\begin{lem}
\label{SameFixed}
Let  $\Lie G$ be a~semisimple Lie algebra over~$\mathbb{Q}$.
If~$\Lie G$ is $\mathbb{R}$-universal, then $\Lie G$ is inner over some imaginary quadratic extension~$L$
of~$\mathbb{Q}$, and the $*$-action of $\Gal(L/\mathbb{Q})$ is the same as the $*$-action of~$\Gal(\mathbb{C}/\mathbb{R})$.
\end{lem}

\begin{prop}[\protect{\cite[Corollary~3.5, \S~4.2, and Lemma~7.4]{Tits-RepLinIrred}}]\label{Beta=D}
Let~$L$ be the center of the algebra $D_{\Lie G, F}(\lambda):= \mathrm{End}_{\Lie G}({}^F \rho_\lambda)$ $($which,
by Schur's lemma, is a~division algebra$)$.
Then:
\renewcommand{\labelenumi}{{\rm (\theenumi)}}
\begin{enumerate}\itemsep=0pt
\item
\label{Beta=D-invt}
$\Gal(\overline{F}/L) = \{ \sigma \in \Gal(\overline{F}/F) \mid \sigma(\lambda) = \lambda \}$ $($for the $*$-action
of $\Gal(\overline{F}/F))$, and
\item $\beta_{\Lie G, L}(\lambda) = [D_{\Lie G, F}(\lambda)]$, after identifying $H^2(L; \boldsymbol\mu)$ with the
Brauer group of~$L$.
\end{enumerate}
\end{prop}

\begin{proof}
Assume $\Lie G$ is not inner, for otherwise the desired conclusions are obvious.
Let~$L$ be the (unique) minimal extension of~$\mathbb{Q}$ over which $\Lie G$ becomes inner, and choose some dominant
weight~$\lambda$ that is not f\/ixed by any nontrivial element of~$\Gal(L/\mathbb{Q})$.
We know that~$L$ is the center of~$D_{\Lie G, \mathbb{Q}}(\lambda)$ (see Proposition~\ref{Beta=D}(\ref{Beta=D-invt})) and that $D_{\Lie
G, \mathbb{Q}}(\lambda) \otimes_\mathbb{Q} \mathbb{R}$ is a~division algebra (see Lemma~\ref{RunivIff}(\ref{RunivIff-DivAlg})).
Since~$L$ is not $\mathbb{Q}$, this implies it is an imaginary quadratic extension of~$\mathbb{Q}$.

It now suf\/f\/ices to show that every every $\Gal(\mathbb{C}/\mathbb{R})$-invariant dominant weight~$\lambda$ is also
$\Gal(L/\mathbb{Q})$-invariant.
Suppose not.
By replacing~$\lambda$ with an appropriate positive integer multiple, we may assume it is in the root lattice, so
$D_{\Lie G,\mathbb{R}}(\lambda) = \mathbb{R}$ (see Remark~\ref{OnlyZStar}).
Therefore, we must have $D_{\Lie G,\mathbb{Q}}(\lambda) = \mathbb{Q}$ (see Lemma~\ref{RunivIff}(\ref{RunivIff-D})).
However, this contradicts Proposition~\ref{Beta=D}(\ref{Beta=D-invt}), since~$\lambda$ is not
$\Gal(\overline{\mathbb{Q}}/\mathbb{Q})$-invariant.
\end{proof}

% Bibliography excerpt retained for reader excerpts of
%   Dave Witte Morris, "A Cohomological Proof that Real Representations of
%   Semisimple Lie Algebras Have Q-Forms", SIGMA 11 (2015) 034,
%   DOI 10.3842/SIGMA.2015.034, CC BY 4.0.
%
% Every entry below is copied character for character from the author's own
% thebibliography environment of the pinned publisher source
% (pinned-original-source.tex lines 954-1016).  No entry was rewritten,
% reordered, merged, split or invented, and no citation key was changed.
\begin{thebibliography}{99}

\bibitem{Morris-RRepsQForms}
Morris D., Real representations of semisimple {L}ie algebras have {${\mathbb
  Q}$}-forms, in Algebraic Groups and Arithmetic, Tata Institute of Fundamental
  Research, Mumbai, 2004, 469--490.


\bibitem{PrasadRapinchuk-Prescribed}
Prasad G., Rapinchuk A., On the existence of isotropic forms of semi-simple
  algebraic groups over number f\/ields with prescribed local behavior,
  \href{http://dx.doi.org/10.1016/j.aim.2006.01.001}{\textit{Adv. Math.}} \textbf{207} (2006), 646--660.


\bibitem{Tits-RepLinIrred}
Tits J., Repr\'esentations lin\'eaires irr\'eductibles d'un groupe r\'eductif
  sur un corps quelconque, \href{http://dx.doi.org/10.1515/crll.1971.247.196}{\textit{J.~Reine Angew. Math.}} \textbf{247} (1971),
  196--220.



\end{thebibliography}
\end{document}
